\section{Implementation Details}
\label{sec:implementation}

\subsection{Technology Stack}

\begin{table}[H]
    \centering
    \caption{Technology Stack}
    \begin{tabular}{@{}lll@{}}
        \toprule
        \textbf{Layer} & \textbf{Technology} & \textbf{Version} \\ \midrule
        Frontend & React & 19 \\
        Frontend Language & TypeScript & 5.x \\
        UI Framework & Ant Design & 6 \\
        Graph Visualization & ReactFlow & latest \\
        Video Player & @mux/mux-video-react & latest \\
        Backend & Python & 3.11 \\
        Web Framework & Django & 5.2 \\
        API Framework & Django REST Framework & 3.15 \\
        WSGI Server & Gunicorn & 21.2+ \\
        Agent Framework & LangGraph & latest \\
        ASR Service & Alibaba DashScope Qwen3-ASR & --- \\
        LLM & Qwen2.5 / Qwen3-max & --- \\
        Vector DB & ChromaDB & 0.4.x \\
        Embeddings & sentence-transformers & 2.x \\
        Video Processing & FFmpeg & 6.x \\
        Computer Vision & OpenCV & 4.x \\
        Cloud Storage & Tencent COS & --- \\
        Reverse Proxy & nginx & 1.27 \\
        Containerisation & Docker Compose & v2 \\
        Database & SQLite (dev/prod-small) & --- \\ \bottomrule
    \end{tabular}
\end{table}

\subsection{Docker Deployment Architecture}

The system is containerised with Docker Compose using three services and a shared named volume.

\subsubsection{Container Layout}

\begin{itemize}
    \item \textbf{web}: Django application served by Gunicorn (4 workers, timeout 120s). Handles all REST API requests and SSE streaming.
    \item \textbf{worker}: Identical image to \texttt{web}, started with \texttt{SERVICE=worker} environment variable. Runs \texttt{python manage.py process\_async\_task} for the video processing pipeline.
    \item \textbf{frontend}: Two-stage image --- Node.js builds the React bundle, nginx serves the static files with SPA routing.
\end{itemize}

The \texttt{web} and \texttt{worker} services share a named Docker volume (\texttt{lecturemind\_data}) mounted at \texttt{/data}, providing a common location for SQLite database, media files, ChromaDB index, and application logs.

\subsubsection{Runtime Environment Injection}

React applications built with Create React App bake environment variables into the bundle at build time, making it impossible to change the backend URL without rebuilding. We solve this by writing a small JavaScript file at container start:

\begin{lstlisting}[language=text, caption={nginx Entrypoint: Runtime Config Injection}]
# docker-entrypoint.sh (frontend)
cat > /usr/share/nginx/html/env-config.js <<EOF
window.__ENV__ = { API_PREFIX: "${API_PREFIX:-http://localhost:8000}" };
EOF
exec nginx -g 'daemon off;'
\end{lstlisting}

The \texttt{index.html} loads this file as a \texttt{<script>} tag, and \texttt{src/config.ts} reads the value:

\begin{lstlisting}[language=text, caption={Frontend Runtime Config (config.ts)}]
export const API_PREFIX: string =
    window.__ENV__?.API_PREFIX ?? "http://127.0.0.1:8000"
\end{lstlisting}

This allows changing the backend URL purely through the \texttt{API\_PREFIX} environment variable with no image rebuild.

\subsubsection{Configuration via Environment Variables}

All paths and ports are configurable through a \texttt{.env} file:

\begin{table}[H]
    \centering
    \caption{Key Environment Variables}
    \begin{tabular}{@{}lll@{}}
        \toprule
        \textbf{Variable} & \textbf{Default} & \textbf{Description} \\ \midrule
        \texttt{BACKEND\_PORT}        & 8000                      & Gunicorn/runserver port \\
        \texttt{FRONTEND\_PORT}       & 3000                      & nginx port \\
        \texttt{MEDIA\_ROOT}          & \texttt{\textless{}BASE\_DIR\textgreater{}/media}  & All uploaded/generated files \\
        \texttt{MEDIA\_AUDIO\_DIR}    & \texttt{\textless{}MEDIA\_ROOT\textgreater{}/audio}      & WAV files for ASR \\
        \texttt{MEDIA\_STREAMS\_DIR}  & \texttt{\textless{}MEDIA\_ROOT\textgreater{}/streams}    & HLS segments \\
        \texttt{MEDIA\_THUMBNAILS\_DIR} & \texttt{\textless{}MEDIA\_ROOT\textgreater{}/thumbnails} & Slide thumbnails \\
        \texttt{CHROMA\_PERSIST\_DIR} & \texttt{\textless{}MEDIA\_ROOT\textgreater{}/chromadb}   & Vector DB storage \\
        \texttt{DB\_PATH}             & \texttt{\textless{}BASE\_DIR\textgreater{}/db.sqlite3}   & SQLite database path \\
        \texttt{LOG\_DIR}             & \texttt{\textless{}BASE\_DIR\textgreater{}/logs}         & Application logs \\ \bottomrule
    \end{tabular}
\end{table}

\subsection{Dual-Resolution Thumbnail Pipeline}

To improve OCR quality, thumbnail generation produces two image sizes per detected slide transition:

\begin{enumerate}
    \item \textbf{200px width} (\texttt{image} field) --- stored for fast web display and slide navigation UI
    \item \textbf{1920px width} (\texttt{image\_high\_res} field) --- stored for high-fidelity OCR processing
\end{enumerate}

The Slide OCR task reads \texttt{image\_high\_res} when available and falls back to \texttt{image}. Using full-resolution thumbnails dramatically improves OCR accuracy for small-font text, equations, and code listings commonly found in lecture slides.

\subsection{API Endpoints}

Key REST API endpoints:

\begin{table}[H]
    \centering
    \caption{Core API Endpoints}
    \begin{tabular}{@{}lll@{}}
        \toprule
        \textbf{Method} & \textbf{Endpoint} & \textbf{Description} \\ \midrule
        POST & /api/videos/upload/ & Upload video \\
        POST & /api/videos/process/ & Trigger processing \\
        GET & /api/videos/\{id\}/transcript/ & Get transcript \\
        GET & /api/videos/\{id\}/sections/ & Get sections \\
        GET & /api/videos/\{id\}/knowledge/ & Get knowledge points \\
        GET & /api/videos/\{id\}/summary/ & Get summary \\
        GET & /api/videos/\{id\}/mindmap/ & Get mindmap \\
        POST & /api/videos/\{id\}/chat/stream/ & Fast RAG chat (SSE) \\
        POST & /api/videos/\{id\}/agent/stream/ & Agentic RAG chat (SSE) \\
        GET & /api/tasks/video/\{id\}/ & Get task status \\
        GET & /api/config/ & List system configuration \\
        POST & /api/config/update/ & Update configuration \\
        GET & /api/health/ & Health check \\ \bottomrule
    \end{tabular}
\end{table}

\subsection{LLM Integration}

The LLM client provides an OpenAI-compatible abstraction:

\begin{lstlisting}[language=Python, caption={LLM Client Interface}]
class LLMClient:
    def __init__(self, api_key, base_url, model):
        self.client = OpenAI(api_key=api_key, base_url=base_url)
        self.model = model
    
    def chat(self, messages, temperature=0.5, max_tokens=2048):
        response = self.client.chat.completions.create(
            model=self.model,
            messages=messages,
            temperature=temperature,
            max_tokens=max_tokens
        )
        return response.choices[0].message.content
    
    def stream_chat(self, messages, tools=None):
        # Generator yielding tokens as they arrive
        ...
    
    def chat_with_tools(self, messages, tools, temperature=0.3):
        # Function calling support
        ...
\end{lstlisting}

This abstraction enables easy model swapping and consistent interface across different LLM providers.

\subsection{Error Handling and Logging}

Comprehensive error handling ensures system reliability:

\begin{itemize}
    \item \textbf{Task-level errors:} Caught and recorded in AsyncTaskItem.error\_message
    \item \textbf{Cascade failures:} Downstream tasks automatically marked as error
    \item \textbf{LLM call logging:} All requests/responses saved to logs/llm\_calls/
    \item \textbf{Retry mechanism:} Failed tasks can be reset and retried via API
    \item \textbf{RAG fallback chain:} Agentic RAG $\to$ Fast RAG $\to$ LLM Direct on failure
\end{itemize}

\subsection{Memory Optimization}

Several techniques reduce memory usage for deployment on modest hardware:

\begin{itemize}
    \item Lazy initialization of embedding model (loaded on first use)
    \item Semantic similarity check disabled by default (saves 200MB)
    \item Frame preprocessing reduces SSIM memory to ~1MB
    \item Batch processing for embedding upserts (100 items per batch)
    \item Streaming LLM responses avoid buffering full responses
\end{itemize}

These optimizations enable the system to run on machines with 8GB RAM.



